\documentclass[10pt,conference]{IEEEtran}

% Packages
\usepackage{cite}
\usepackage{amsmath,amssymb,amsfonts}
\usepackage{algorithmic}
\usepackage{graphicx}
\usepackage{textcomp}
\usepackage{xcolor}
\usepackage{booktabs}
\usepackage{hyperref}

% Correct bad hyphenation here
\hyphenation{op-tical net-works semi-conduc-tor}

\begin{document}

% Paper title
\title{Title of Your Research Paper}

% Author names and affiliations
\author{
\IEEEauthorblockN{Author Name}
\IEEEauthorblockA{
    Department Name\\
    Institution Name\\
    City, Country\\
    email@institution.edu}
\and
\IEEEauthorblockN{Second Author}
\IEEEauthorblockA{
    Department Name\\
    Institution Name\\
    City, Country\\
    email@institution.edu}
}

% Make the title area
\maketitle

% Abstract
\begin{abstract}
This is the abstract of your paper. It should be 150-250 words and summarize the 
purpose, methods, results, and conclusions of your research. The abstract should 
be self-contained and readable without the full paper. Include specific metrics 
and quantitative results where possible.
\end{abstract}

% Keywords
\begin{IEEEkeywords}
keyword1, keyword2, keyword3, keyword4, keyword5
\end{IEEEkeywords}

\IEEEpeerreviewmaketitle

% Section 1: Introduction
\section{Introduction}
\label{sec:introduction}

Start with a compelling real-world problem or motivation. Build from general to 
specific context. Clearly state the research gap and your contributions.

The main contributions of this paper are:
\begin{itemize}
    \item Contribution 1: Description of your first contribution.
    \item Contribution 2: Description of your second contribution.
    \item Contribution 3: Description of your third contribution.
\end{itemize}

The remainder of this paper is organized as follows. Section \ref{sec:related} 
reviews related work. Section \ref{sec:method} describes the proposed method. 
Section \ref{sec:experiments} presents experimental results. Section 
\ref{sec:conclusion} concludes the paper.

% Section 2: Related Work
\section{Related Work}
\label{sec:related}

Review relevant literature and position your work. Group related work by theme 
or approach. Compare explicitly: "Unlike [X] which focuses on Y, our approach 
addresses Z..."

% Section 3: Methodology
\section{Methodology}
\label{sec:method}

Describe your proposed method in detail. Include architecture diagrams, 
algorithms, and design rationale.

\subsection{Architecture Overview}

Provide a high-level description of your system or method.

\subsection{Key Components}

Describe each component in detail.

\subsection{Algorithm}

Present key algorithms using pseudocode:

\begin{algorithm}
\caption{Algorithm Name}
\begin{algorithmic}[1]
\REQUIRE Input parameters
\ENSURE Output results
\STATE Step 1
\STATE Step 2
\STATE Step 3
\RETURN result
\end{algorithmic}
\end{algorithm}

% Section 4: Experiments
\section{Experiments}
\label{sec:experiments}

\subsection{Experimental Setup}

Describe datasets, evaluation metrics, baselines, and implementation details.

\subsection{Results}

Present quantitative results in tables:

\begin{table}[htbp]
\caption{Performance Comparison}
\label{tab:results}
\centering
\begin{tabular}{@{}llll@{}}
\toprule
Method & Metric 1 & Metric 2 & Metric 3 \\
\midrule
Baseline & 0.80 & 0.75 & 0.82 \\
Proposed & \textbf{0.91} & \textbf{0.89} & \textbf{0.92} \\
\bottomrule
\end{tabular}
\end{table}

\subsection{Analysis}

Provide ablation studies, error analysis, and qualitative examples.

% Section 5: Conclusion
\section{Conclusion}
\label{sec:conclusion}

Summarize your contributions and their significance. Discuss limitations and 
future work.

% Acknowledgments (optional)
% \section*{Acknowledgments}
% Thank funding sources, collaborators, etc.

% References
\bibliographystyle{IEEEtran}
% \bibliography{references}

% Example manual references (or use .bib file)
\begin{thebibliography}{1}

\bibitem{example1}
A. Author and B. Author, ``Title of paper,'' in \emph{Proc. Conference Name}, 
Year, pp. 123--456.

\bibitem{example2}
C. Author, ``Title of journal article,'' \emph{Journal Name}, vol. X, no. Y, 
pp. 123--456, Month Year.

\end{thebibliography}

\end{document}
